\section{Introduction}

\begin{figure}[t]
    \centering
    \includegraphics[width=0.92\linewidth]{figures/gradcam_results_2x2.png}
    \caption{Illustrating the targeted localization feature ($P_{t,c}$; Eq.~\ref{eq:gradcam_pm}) within the ASCG framework. It shows the feasibility of spatially targeted interventions, i.e., \emph{where} to steer (\S\ref{subsec:targeted}); notably, harmful attributes are identified across denoising timesteps as well, facilitating temporal selective intervention, i.e., \emph{when} to steer (\S\ref{subsec:selective}).}
    \label{figure/gradcam}
\end{figure}

The rapid evolution of \emph{text-to-image (T2I)} diffusion models has enabled the generation of high-fidelity images with user-specified semantics~\citep{rombach2022high,ramesh2022hierarchical}.
As these models become increasingly practical and integrated into diverse applications, their ability to synthesize visual concepts from simple textual prompts raises significant concerns regarding the generation of harmful or inappropriate content, such as explicit nudity or copyrighted material \citep{rando2022red,schramowski2023safe,qu2023unsafe}.
To mitigate these risks, a growing body of research has established various safety mechanisms, which can be broadly categorized into: \emph{1) Training-based approaches}, which involve fine-tuning or unlearning methods to permanently modify model parameters and remove harmful concepts~\citep{gandikota2023erasing,kim2023towards}; and \emph{2) Training-free inference-time guidance}, which steers the denoising process away from unsafe regions using negative prompts or external safety classifiers~\citep{schramowski2023safe,brack2023sega}.

However, these methods often involve a trade-off between safety and generative integrity due to their inflexible or indiscriminate nature: \emph{1) Training-based approaches} typically offer limited adaptability; because the model parameters are fundamentally altered, the intervention remains static regardless of the generation context, potentially leading to unintended degradation in image quality.
\emph{2) Training-free inference-time guidance} often applies suppression in a globalized manner across all timesteps and spatial locations; by applying uniform suppression regardless of actual necessity, existing methods frequently over-correct benign regions, inducing unnecessary semantic drift even when safety risks are localized or entirely absent~\citep{gandikota2023erasing,schramowski2023safe}.
To address these challenges, we propose \emph{Adaptive Spatial Classifier Guidance (ASCG)}, a necessity-driven framework that employs probabilistic monitoring to activate guidance only upon detecting potential harm.
\begin{itemize}
    \item Figure \ref{figure/gradcam} illustrates the core mechanism for steering the generation of T2I models within the ASCG framework.
    \item ASCG enables surgical inference-time guidance for specific latent features by leveraging the spatial localization capabilities of Grad-CAM~\citep{selvaraju2017grad}, effectively suppressing harmful concepts while minimizing collateral damage to benign components during the generation process.
    \item At each denoising step, a Grad-CAM operator yields a harmful-class saliency whose global score gates guidance intervention (i.e., \emph{when} to steer; \S\ref{subsec:selective}) and whose spatial map defines a binary mask that localizes harmful regions (i.e., \emph{where} to steer; \S\ref{subsec:targeted}).
    \item We not only empirically validate the efficacy of ASCG on nudity erasure tasks but also address the inherent limitations of existing benchmarks by incorporating a newly proposed VLM-based evaluation framework that ensures a more robust assessment (\S\ref{sec:exp}).
\end{itemize}

% Text-to-image diffusion models have achieved strong image quality and prompt alignment. At the same time, they raise persistent concerns regarding the generation of unsafe or inappropriate content, such as nudity or violence. Addressing these concerns requires suppressing harmful content while preserving the semantic intent of the prompt and the model’s general generative capability.

% To address this challenge, a growing body of work has proposed safety mechanisms for diffusion models. One line of research focuses on \emph{training-based} approaches, including fine-tuning and unlearning methods that permanently modify model parameters to remove harmful concepts \textbf{(e.g., ESD; SDD)}. While such methods can enforce strong suppression, they require substantial training resources and lack adaptability: once the model is modified, its behavior is affected even when no harmful content is present, often leading to degraded image quality or reduced generation flexibility.

% An alternative direction explores \emph{training-free inference-time guidance}, most notably by manipulating classifier-free guidance (CFG) to suppress target concepts. Methods such as SLD and SEGA subtract concept-specific guidance directions during denoising, applying suppression globally across all timesteps and spatial locations. Although effective at reducing explicit activation of undesired concepts, these approaches implicitly assume that harmful content is present whenever a target concept is specified. As a result, guidance is applied uniformly, frequently overcorrecting benign regions and inducing visual artifacts or semantic drift from the original prompt.

% Prompt-level modification and filtering introduce a related limitation. By altering or removing prompt tokens associated with unsafe concepts, these methods achieve safety by redefining the input itself, inevitably compromising the user’s original semantic intent. Even when model parameters remain unchanged, such hard filtering can introduce distribution shifts that degrade generation quality.

% Across both training-based and training-free approaches, a shared limitation emerges: safety is enforced through \emph{constant, global intervention}, regardless of whether harmful content actually arises during generation or where it appears in the image. In practice, unsafe content is often sparse, implicit, or localized, and may not appear at all for many prompts. This mismatch leads to unnecessary suppression of benign content and highlights the need for safety mechanisms that are selective rather than globally applied.

% These observations motivate a different perspective on diffusion safety. Effective safety intervention should be \emph{conditional}, activating only when harmful content emerges, and \emph{spatially selective}, targeting only the regions responsible for unsafe signals. Moreover, such intervention should avoid permanent modification of model parameters, preserving the model’s general generative capability.

% In this work, we explore this perspective through \textbf{Adaptive Spatial Classifier Guidance (ASCG)}, an inference-time framework that enforces safety without modifying diffusion model parameters. ASCG employs a monitoring algorithm to detect harmful signals in diffusion latents and, when triggered, applies Grad-CAM--based spatial guidance to selectively suppress harmful regions. While ASCG is broadly applicable to various safety settings, this workshop paper focuses on \emph{nudity removal} as a representative case study.

% Our contributions are summarized as follows:
% \begin{itemize}
%     \item We propose \textbf{Adaptive Spatial Classifier Guidance (ASCG)}, a training-free inference-time safety framework that enforces safety through conditional and spatially selective guidance, without modifying diffusion model parameters or altering input prompts.
%     \item We show that existing global guidance and prompt-level suppression methods are inherently misaligned with semantic preservation, and that selective intervention provides a more principled trade-off between safety and \emph{semantic integrity}.
%     \item Through nudity removal experiments evaluated with vision--language models, we demonstrate that ASCG improves perceptual safety while maintaining semantic integrity compared to globally applied guidance methods.
% \end{itemize}

%%%%%%%%%%%%%%%%%%%%%%%%%%%%%%%%%%%%%%%%%%%%%%%%%%%%%%%%%%%%%%%%%%%%%%%%%%%%%%%%
% \section{Related Work}

% \paragraph{Safety in T2I Models.} Building on the growing body of work exposing vulnerabilities in T2I models~\citep{}, existing safety mechanisms generally follow two main routes: model editing and inference-time guidance.

% Model editing methods permanently alter parameters to suppress harmful concepts, ranging from iterative fine-tuning like ESD~\citep{gandikota2023erasing} and SDD~\citep{kim2023towards} to closed-form updates such as UCE~\citep{gandikota2024unified} and RECE~\citep{gong2024reliable}. While effective, these methods suffer from irreversible weight modifications, often degrading general image quality and limiting flexibility for benign prompts.

% In contrast, inference-time guidance intervenes dynamically during image generation without modifying model weights. Early works like SLD~\citep{schramowski2023safe} and SEGA~\citep{brack2023sega} steer the denoising trajectory across the entire latent image, inevitably causing unintended shifts in global structure or background due to their uniform application. While SAFREE~\citep{yoon2024safree} refines this by projecting harmful tokens orthogonal to the toxic subspace, such modification of input embeddings risks semantic inconsistency and remains vulnerable to implicit adversarial prompts lacking explicit toxic keywords.

% By fundamentally altering the direction of core concepts within the text space, it often disrupts the overall coherence of the generated image.

% \paragraph{Safety in Text-to-Image Models.} A growing body of work has explored safety mechanisms for text-to-image diffusion models, aiming to suppress harmful concepts while preserving semantic intent and generative capability. Rather than categorizing prior work solely as training-based or training-free, we organize existing approaches according to whether they modify the internal parameters of the diffusion model. This distinction is fundamental: while parameter modification risks irreversible interference with benign concepts due to the entanglement of latent representations, parameter-unmodified methods provide context-dependent intervention that preserves the model's original generative prior for safe prompts.

% \paragraph{Model Editing Methods.} A large class of safety approaches enforces suppression by directly altering model weights. This category is primarily composed of \textbf{training-based methods} that remove or weaken harmful concepts through iterative fine-tuning. Representative examples include direct concept removal (ESD)~\citep{gandikota2023erasing} and SDD~\citep{kim2023towards}, the latter of which utilizes self-distillation to guide model weights toward a safer distribution. While effective at targeted suppression, these methods permanently alter the diffusion model, often leading to unintended degradation of image quality and reduced generalization for benign prompts.

% Importantly, parameter modification is not limited to iterative training. Recent works explore training-free weight editing, such as UCE \citep{gandikota2024unified} and RECE~\citep{gong2024reliable}. Although these methods utilize closed-form updates to eliminate the computational overhead of fine-tuning, they still irreversibly modify the model’s weight matrix. This permanent alteration makes them less flexible for adaptive safety control, as the model remains in a constrained state even when generating entirely benign content.

% \paragraph{Inference-time Guidance.} In contrast, parameter-unmodified approaches preserve the original diffusion model and intervene only at inference time. Early methods in this category, such as Safe Latent Diffusion (SLD)~\citep{schramowski2023safe} and Semantic Guidance (SEGA)~\citep{brack2023sega}, manipulate classifier-free guidance (CFG) to steer trajectories away from harmful concepts. These methods apply suppression globally across all spatial locations. More recent work, SAFREE~\citep{yoon2024safree}, introduces a stronger safeguard by jointly operating in text and latent spaces. However, SAFREE primarily relies on token-level and latent-level interventions that are applied globally to the image as a whole, without explicitly reasoning about where harmful content emerges in the spatial plane.

% \paragraph{Classifier Guidance and Trajectory-Based Control.} Classifier guidance modifies the denoising trajectory of diffusion models by incorporating gradients from auxiliary classifiers~\citep{Dhariwal2021Diffusion, nichol2022glide}. Beyond globally applied guidance, trajectory-based methods such as sparse repellency (SR)~\citep{kirchhof2024shielded} penalize diffusion trajectories only when the expected output approaches a protected reference set. While SR effectively reduces unnecessary intervention at the sample level through sparse triggers that are inactive most of the time, its repellency mechanism relies on global distance metrics within the latent space. Consequently, it lacks spatial granularity, making it difficult to suppress localized harmful features without affecting the global semantic structure of the generated image.

% \textbf{Adaptive Spatial Classifier Guidance (ASCG)} belongs to the class of parameter-unmodified safety mechanisms but departs fundamentally from prior global guidance approaches. Unlike SLD~\citep{schramowski2023safe}, which applies guidance uniformly, and SAFREE~\citep{yoon2024safree}, which focuses on global latent modulation, ASCG explicitly introduces \textbf{spatial adaptivity}.

% ASCG determines when to intervene through probabilistic monitoring of harmful signal emergence and where to intervene via \textbf{Grad-CAM--based spatial localization}. By leveraging CDF-normalized attribution maps, ASCG targets specific regions where harmful concepts emerge while preserving benign content elsewhere. This spatially selective design directly addresses the overcorrection and semantic drift observed in prior methods, aligning safety intervention more closely with the localized nature of harmful visual content.

%%%%%%%%%%%%%%%%%%%%%%%%%%%%%%%%%%%%%%%%%%%%%%%%%%%%%%%%%%%%%%%%%%%%%%%%%%%%%%%%
